%!TEX root = ../../thesis.tex

\subsection{Branches}\label{sec:target-branches}

% TODO: The branching-maze target, and the strongest argument in the thesis for
% planning rather than greedy value maximization --- give it that framing.
%
% Structure: the input space is a DAG of rooms rather than one linear passcode. Each
% room opens a few distinct-letter doors, wrong letters kill the process, and several
% leaves are goals of differing value. The layout stresses several capabilities at
% once, and each deserves a sentence:
%   - the global optimum lies several decision-heavy steps deep,
%   - a trivially repeatable "just survive" corridor competes with it,
%   - a shallow but well-rewarded sled is an early local optimum,
%   - goals branch off the correct deep path, so a greedy climber stops short,
%   - a merge room is reachable by two distinct prefixes, testing generalization
%     across convergent paths,
%   - dead ends sit directly on the deep branch.
%
% Reward mirrors sequence (\cref{sec:target-sequence}): a living reward per surviving
% step and a terminal payout per distinct block, sized so that the goal sleds rank the
% goals. Include the maze as a figure --- this target is unreadable in prose alone.
%
% State the claim it is built to test: a greedy or coverage-guided climber takes the
% sled or the corridor, while a planner with a learned model should look past both.
% If the planner does not beat the climber here, the planning claim of RQ3
% (\cref{sec:introduction-rqs}) fails on the target designed to make it easiest ---
% which is why this result is reported whichever way it comes out.
